\setlength{\absparsep}{12pt}
\begin{resumo}
Por que os epiciclos de Ptolomeu e Copérnico descreviam tão bem o céu? Desenvolvemos um simulador do Sistema Solar que integra a gravitação mútua do Sol e dos oito planetas a partir de condições iniciais do JPL Horizons, com o integrador simplético de Yoshida de quarta ordem e passo fixo de meio dia. Contra a efeméride DE440s, numa varredura de $-100$ a $+100$ anos, o desvio fica abaixo de 0,5\,\% da distância ao Sol; no pior caso, Mercúrio, chega a 0,46\,\% e a 13 minutos de arco na direção vista da Terra. A energia oscila sem crescer. Sobre essas órbitas, calculamos a série de Fourier complexa parametrizada no tempo, com forma fechada em funções de Bessel. Até primeira ordem na excentricidade, ela tem a estrutura da construção do \emph{Commentariolus}: um deslocamento fixo de $\tfrac{3}{2}ae$, um deferente de raio $a$ e um epiciclo de raio $\tfrac{1}{2}ae$ com o dobro da frequência. Os raios de Copérnico para Marte, Júpiter e Saturno concordam com esses valores em até 8\,\%, e a razão de três para um entre os seus epiciclos, herdada da substituição do equante, coincide com a da série. O equante de Ptolomeu, usado como controle, tem o mesmo conteúdo de primeira ordem: com parâmetros modernos e o deferente de Copérnico igual a $|c_1|$, ele, o modelo de Copérnico e a série erram a posição de Marte em cerca de 0,007\,UA, os três em ordem $e^2$. A série explica por que os dois modelos funcionam.

\textbf{Palavras-chave}: epiciclos. séries de Fourier. problema de $N$ corpos. integrador simplético. Copérnico.
\end{resumo}
